% Dibangkitkan oleh revisi/hasil/gerbang_batch/bab5_jalan.py. Jangan diubah dengan tangan.
\begin{table}[htbp]
\caption{Pelanggaran keras pemeriksa jadwal pada data uji (I-2)}
\label{tab:uji-layak}
\centering\footnotesize
\setlength{\tabcolsep}{4pt}
\begin{adjustbox}{max width=\textwidth}
\begin{tabular}{lrrr}
\toprule
Metode & Episode & \shortstack[r]{Pelanggaran\\keras} & \shortstack[r]{Episode tanpa\\pemeriksa} \\
\midrule
ARLS-PPO & 33.210 & 0 & 0 \\
AOR & 11.070 & 0 & 0 \\
Random + penjaga & 11.070 & 0 & 0 \\
Fixed-L3 + penjaga & 11.070 & 0 & 0 \\
Boosting-IL + penjaga & 11.070 & 0 & 0 \\
ARLS-BC & 33.210 & 0 & 0 \\
ARLS-PPO tanpa penjaga & 33.210 & 0 & 0 \\
Sequential Lookahead & 11.070 & 0 & 0 \\
\bottomrule
\end{tabular}
\end{adjustbox}
\par\smallskip\footnotesize\textit{Catatan:} Episode metode PPO dihitung untuk tiga \emph{seed} model.
\end{table}
